\chapter{Expressions and Equations}

All of algebra is built on the idea of \textit{expressions} and \textit{equations}.
Expressions and equations allow us to express ideas compactly using symbols, and they will continue to be critical for the rest of your mathematical career.

\begin{definition}
Expression: An expression is a group of symbols and operators that represent a value. An expression never includes an equals sign.
\end{definition}

\begin{definition}
Equation: An equation consists of two expressions related by equality, in the form $a=b$ where $a$ and $b$ are expressions.
\end{definition}

\noindent Here are a few examples of expressions:
\[1 + 1\]
\[x + 3\]
\[3y - 17\]
\[\tau - 1\]

\noindent Here are some equations:
\[A=\frac{1}{2}\tau r^{2}\]
\[x+3=4\]
\[x=1\]
\[1+1=2\]
\[\tau=6.28318530717\dots\]

Expressions and equations can contain constants, variables, or a mix of the two. 
A variable is a quantity, often represented by a letter, whose value may change.
A constant is a fixed value.
Examples of constants include $1$, $6$, $\tau$, and $12.463$.
\begin{definition}
    \textit{Numerical Expression}: A numerical expression contains only constants
\end{definition}
\begin{definition}
    \textit{Algebraic Expression}: An expression containing at least one variable.
\end{definition}
When one is given an expression, their first goal should be to simplify the expression.
An expression is considered simplified when all possible operations have been performed.
For example, if one is given the expression $1+1$, they will usually simplify it to $2$.
Algebraic expressions may be simplified as well.
In this case, one must apply a technique called "combining like terms".
Often variables will have numbers in front of them, these numbers are called coefficients.
An example of a coefficient is the "3.5" in $3.5x$.
This expression is equal to $3.5 * x$.
When you combine like terms, you first group together all the variables of the same type.
Then, you add up their coefficients, and multiply the variable by that coefficient.
This new expression can replace all instances of that variable in the original expression.
\newline \newline Example:
\[x + 3y + 2x - 0.5y\]
\[= 2x + x + 3y - 0.5y\]
\[= 3x + 2.5y\]
Note that $x$ is the same as $1x$.
\newline \newline
\indent As you saw in definition 2, an equation consists of two expressions separated by an equality.
Therefore, if one is given an equation, they may simplify each side, and those sides will remain equal.
Example:
\[x + 2x + 1 = 3 - y + 4 + y\]
\[3x + 1 = 7\]
\indent When a single variable appears in an equation, it is called an unknown.
This is because that variable can take on only 1 value.
In the equation $3x+1=7$, $x=2$.
This can be verified by substituting $2$ for $x$:
\[3(2)+1=7\]
\[6+1=7\]
\[7=7\]
Which is true.
\newline \indent Because $x$ can only take on one value, it is called an "unknown" until its value is determined, then one may refer to it as a constant.
Much of algebra focuses on determining the values of unknowns, which we will soon learn.
\newline \indent In an equation like $y + 1 = 3x -7$, $x$ may take on many values (infinitely many in fact), and so may $y$.
That is why, in this context, we refer to $x$ and $y$ and \textit{variables}.
\Large{Exercises:}
1. Determine whether each of the following is an equation or an expression:


